% !PRINTED 43
\weberpage{43}
%<<CONT
\[
S_0 = \int \frac{dx}{\sqrt{f(x)}}, \quad S_2 = \int \frac{x^2\,dx}{\sqrt{f(x)}}, \quad S_{-1} = \int \frac{dx}{(x-\alpha)\sqrt{f(x)}},
\]
while $S_1 = \displaystyle\int \frac{x\,dx}{\sqrt{f(x)}}$ is reduced by the substitution $x^2 = y$ to the following non-elliptic integral:
\[
\frac{1}{2}\int \frac{dy}{\sqrt{y^2 + a_2 y + a_4}}
\]
Here $S_0$ is of the first, $S_2$ of the second, $S_{-1}$ of the third kind. For the \gesperrt{Legendre} normal form we put $f(x) = (1-x^2)(1-\varkappa^2 x^2)$ and take as the normal integral of the second kind not $S_2$ itself, but $S_0 - \varkappa^2 S_2$. This gives the normal integrals of the three kinds:
\[
\int \frac{dx}{\sqrt{(1-x^2)(1-\varkappa^2 x^2)}}, \quad \int \sqrt{\frac{1-\varkappa^2 x^2}{1-x^2}}\,dx, \quad \int \frac{dx}{(x-\alpha)\sqrt{(1-x^2)(1-\varkappa^2 x^2)}}
\]
and, if one puts $x = \sin\varphi$:
\[
\int \frac{d\varphi}{\sqrt{1-\varkappa^2\sin^2\varphi}}, \quad \int \sqrt{1-\varkappa^2\sin^2\varphi}\,d\varphi, \quad \int \frac{d\varphi}{(\sin\varphi - \alpha)\sqrt{1-\varkappa^2\sin^2\varphi}}.
\]

The same is obtained if one takes the elliptic differential in the normal form
\[
\frac{dz}{\sqrt{z(1-z)(1-\varkappa^2 z)}}
\]
and proceeds according to Case II.

\begin{center}
\textbf{\S\ 13. The \gesperrt{Addition Theorem.}}
\end{center}

The \gesperrt{addition theorem} of the elliptic integrals, discovered by \gesperrt{Euler}, consists in the theorem, of fundamental importance for the whole of the theory that follows, that if, of the three pairs of values
\[
x_1, \sqrt{f(x_1)}; \quad x_2, \sqrt{f(x_2)}; \quad x_3, \sqrt{f(x_3)}
\]
two are assumed to be given, the third can be determined by \gesperrt{algebraic means} in such a way that the differential equation
\weq{1}{\frac{dx_1}{\sqrt{f(x_1)}} + \frac{dx_2}{\sqrt{f(x_2)}} + \frac{dx_3}{\sqrt{f(x_3)}} = 0}
is satisfied. In this, $f(x)$ is a given function of the third or fourth degree. This is a special case of the great
%CONT>>
